\section*{Bab \ref{ch:g5:air-pressure-weather} --- Tekanan Udara dan Cuaca}

\begin{solution}{exo:g5:air-pressure-weather:1}
Bola yang sama ditimbang dalam keadaan keras lalu dalam keadaan
kempis: hanya \omterm{def:g2:air-around-us:air}{udara} yang meninggalkannya, dan \omterm{def:g3:levers-and-scales:scale}{neracanya} menunjukkan
lebih sedikit --- selisihnya adalah \omterm{def:g4:weight-and-mass:weight}{berat} \omterm{def:g2:air-around-us:air}{udara} yang lolos. \omterm{def:g2:air-around-us:air}{Udara}
adalah zat yang punya \omterm{def:g3:measuring-mass:mass}{massa}.
\end{solution}

\begin{solution}{exo:g5:air-pressure-weather:2}
Dorongan \omterm{def:g2:air-around-us:air}{udara} yang terperah di dasar \omterm{def:g5:air-pressure-weather:pressure}{atmosfer} pada setiap permukaan
yang disentuhnya. Dorongan itu datang dari \omterm{def:g4:weight-and-mass:weight}{berat} samudra \omterm{def:g2:air-around-us:air}{udara} yang
dalam di atas kita, dan ia mendorong dari segala arah --- ke bawah, ke
samping, dan ke atas.
\end{solution}

\begin{solution}{exo:g5:air-pressure-weather:3}
Karena dorongannya datang berimbang: \omterm{def:g2:air-around-us:air}{udara} di sisi lain telapak
tanganmu --- dan di dalam tubuhmu --- mendorong balik sama kuatnya.
Hanya dorongan yang \emph{tidak berimbang}, seperti pada kartunya atau
pada botol yang remuk, yang memperlihatkan kekuatannya.
\end{solution}

\begin{solution}{exo:g5:air-pressure-weather:4}
\omterm{def:g4:weight-and-mass:weight}{Berat} airnya mendorong kartunya ke bawah; \omterm{def:g5:air-pressure-weather:pressure}{atmosfer} mendorong kartunya
ke atas. \omterm{def:g5:air-pressure-weather:pressure}{Atmosfer} menang --- kartunya bertahan dan airnya tertahan di
dalam.
\end{solution}

\begin{solution}{exo:g5:air-pressure-weather:5}
Kamu mengeluarkan \emph{\omterm{def:g2:air-around-us:air}{udara}} dari dalam sedotannya. \omterm{def:g5:air-pressure-weather:pressure}{Atmosfer}, yang
menekan jus di dalam gelas, lalu mendorong jusnya naik lewat sedotan
yang sudah bersih: langitlah yang mengangkatnya.
\end{solution}

\begin{solution}{exo:g5:air-pressure-weather:6}
Menekan mangkuknya rata memerah \omterm{def:g2:air-around-us:air}{udara} di belakangnya keluar. Karena
tidak ada yang mendorong balik dari bawahnya, dorongan \omterm{def:g5:air-pressure-weather:pressure}{atmosfer}
memakukan mangkuknya pada ubin.
\end{solution}

\begin{solution}{exo:g5:air-pressure-weather:7}
Ia melemah --- samudra \omterm{def:g2:air-around-us:air}{udara} yang tersisa di atas untuk menekan ke
bawah menjadi lebih sedikit. Tandanya: telinga berdenging pada
pendakian yang cepat, napas tersengal di \omterm{def:g2:air-around-us:air}{udara} yang tipis (dan teh
yang mendidih terlalu sejuk).
\end{solution}

\begin{solution}{exo:g5:air-pressure-weather:8}
Dorongan \omterm{def:g5:air-pressure-weather:pressure}{atmosfer} --- \omterm{def:g5:air-pressure-weather:pressure}{tekanan udara}. Penurunan yang cepat
memperingatkan badai yang sedang datang: pindahkan pikniknya ke dalam
ruangan.
\end{solution}

\begin{solution}{exo:g5:air-pressure-weather:9}
Ketika disegel di puncak, botolnya menyimpan \omterm{def:g2:air-around-us:air}{udara} gunung yang tipis
di dalamnya; di lembah dorongan dari luar menjadi lebih kuat daripada
dorongan yang terperangkap di dalam, dan botolnya pun melipat. Kalau
dibawa ke arah sebaliknya, botol yang disegel di lembah menyimpan
\omterm{def:g2:air-around-us:air}{udara} yang \emph{kuat} di dalamnya sementara dorongan luarnya melemah
seiring ketinggian: botolnya menggembung tegang --- dan bungkus keripik
yang tersegel melakukan hal yang sama pada setiap perjalanan ke gunung.
\end{solution}

\begin{solution}{exo:g5:air-pressure-weather:10}
\omterm{def:g2:air-around-us:air}{Udara} mengalir tetap keluar dari daerah bertekanan tinggi yang \omterm{def:g4:weight-and-mass:weight}{berat}
menuju daerah bertekanan rendah yang ringan, dan aliran itulah
\omterm{def:g2:air-around-us:wind}{anginnya}. Jadi definisi yang lama melengkapi dirinya sendiri: \omterm{def:g2:air-around-us:wind}{angin}
adalah \omterm{def:g2:air-around-us:air}{udara} yang bergerak --- digerakkan oleh perbedaan \omterm{def:g5:air-pressure-weather:pressure}{tekanan
udara}.
\end{solution}

\begin{solution}{exo:g5:air-pressure-weather:11}
Seiring ketinggian, dorongan \omterm{def:g5:air-pressure-weather:pressure}{atmosfer} pada airnya melemah, sehingga
mendidih menjadi lebih mudah: \omterm{def:g4:changes-of-state:boiling-point}{titik didihnya} meluncur jauh di bawah
\qty{100}{\celsius}. Dan menurut hukum dataran, air yang mendidih
tidak pernah dapat memanjat melewati \omterm{def:g4:changes-of-state:boiling-point}{titik didihnya} --- api tambahan
hanya membuat uap. Jadi air di puncak mendidih dalam keadaan sejuk
\emph{dan terpaku di sana}: tehnya lemah, pastanya keras, fisikanya
puas.
\end{solution}
